\subsection{What a sheaf adds to the complex}
The incidence matrix compares vertex values by ordinary subtraction. A cellular sheaf generalizes the comparison rule. It attaches a vector space, called a \emph{stalk}, to every cell, and specifies linear maps between the spaces on faces and cofaces. To compare values from two vertices, first map them into the common edge space and then subtract. The maps define how the cellwise assignments are compared.

In the present study every stalk is $\R$. Thus all stalk values are scalars and all comparison maps are multiplication by real numbers. The phrase \emph{rank-one sheaf} refers to the dimension of each stalk. A neighborhood with many vertices can still produce a large, high-rank matrix. The more general theory permits different vector-space dimensions and matrix-valued maps, but those possibilities are not needed here \citep{hansen2019}.

The maps are traditionally called restrictions. Our convention uses the face's stalk as the domain and the coface's stalk as the codomain. The displayed arrow $\rho_{\sigma\tau}:\mathcal S(\sigma)\to\mathcal S(\tau)$ specifies the direction. All formulas in this manuscript use this convention consistently.

Why impose compatibility? Suppose a vertex belongs to a triangle through two different intervening edges. The value obtained by mapping the vertex assignment to the triangle must agree along both routes. Otherwise the cancellation in the triangle calculation above can fail. Compatibility guarantees that the two routes have the same map; the incidence signs give them opposite signs. The identity $d_{q+1}d_q=0$ then follows. An arbitrary collection of edge weights can define a positive-semidefinite graph Laplacian without satisfying this stronger higher-dimensional requirement.

It is helpful to separate three ingredients. The complex says which cells exist. The restriction maps say how assignments are compared. The inner products say how large an assignment or a disagreement is. Changing any of these can change a Laplacian. A statement that two constructions have the same cohomology need not imply that they have the same positive spectrum, because cohomology forgets the size information supplied by the inner products.

\subsection{Reading the geometric restriction rule}
The restriction formula multiplies a ratio of positive geometric factors by the labels of newly added vertices. For a vertex the geometric factor is one. For an edge it is its dimensionless length. For a triangle it is the product of its three edge lengths. For an edge contained in a triangle, the restriction divides by the two additional edge lengths and multiplies by the label of the newly added vertex.

Following two successive face relations multiplies two ratios whose intermediate factors cancel. The label products similarly combine into the labels of all newly added vertices. This telescoping is the mechanism behind compatibility. A zero label is safe because it appears only in a product, never in a denominator. Coincident points are different: they would introduce zero geometric denominators and must be handled by the coordinate audit.

We use two choices of vertex labels. All-one labels keep the geometric factors. A zero at the target distinguishes the target algebraically from its neighbors. Both choices use only the target identity and geometry. The sheaf is built before fitting the model, so the eventual response cannot influence its restrictions.
